\section{Limitations}
\label{sec:limitations}

The models are
14M-parameter class, trained 20{,}000 steps from scratch on a synthetic
episode distribution with a GPT-2 tokenizer; nothing here speaks to
pretrained models, natural text, or other scales. The baseline
non-competitiveness holds at the matched budget only: a larger, longer-,
or differently-trained transformer requires its own matched run before any
extension of the claim. The metric is argmax-decoded and
episode-restricted, so by the capacity results of
\citet{nichani2024understanding} it supports no statement about rank,
exact recovery, or how much the 32{,}768-byte state holds; the continuous
diagnostics of Section~\ref{sec:mechanism} are reported only as
bounds. The separation is demonstrated for one fast-weight family (a
delta-rule matrix state) against one vector control and one transformer
configuration; the multi-hop variant does not separate
(Section~\ref{sec:scope}); and the recall ceiling is load-bounded, with
every arm at chance at $K/d = 0.75$. Horizon evaluation extends context
with filler continuation, not with distractor bindings, which is the
easier long-context condition and is disclosed as such. The longest
tested context is 1{,}798 tokens, eight times the binding
span: % <!-- evidence: C3 -->
this paper demonstrates the constant-memory \emph{mechanism} as a
small-scale existence result, not a claim that it holds at the context
lengths (tens of thousands to millions of tokens) where KV-cache cost is
a practical deployment constraint.
